% =========================================================
% doc5-summary.tex — Section 5: Conclusion
% New content (final-rules §1: Section 5, short).
% NOTE: d0-positioning-objectives.md still lacks an objective line for
% this section; suggested objective is stated in the first sentence and
% should be back-ported to d0 (final-rules §6.6).
% =========================================================

\section{Conclusion}\label{sec:conclusion}

We have introduced Polar Semirings, an axiomatization of polarity as a
primitive algebraic notion, and established their basic theory and their
position in the landscape of involutive structures. Three results summarize
the paper. First, the eleven axioms of \Cref{def:ps} are independent and
admit natural models on chains, including Boolean, tropical, and quadratic
models (\Cref{sec:ps}). Second, the axioms decompose into an algebraic
layer, a free reflection layer generated by the involution, and a one-axiom
integration layer, and this architecture yields the Fundamental Structure
Theorem: every Polar Semiring is canonically self-dual, with the polarity
acting as an isomorphism onto the dual structure and the two induced orders
coinciding (\Cref{thm:fundamental}). Third, every bounded Girard algebra
induces a Polar Semiring upon forgetting its residual
(\Cref{thm:bga-implies-ps}), whereas the quadratic model on the unit
interval admits no bounded Girard algebra extension
(\Cref{thm:impossibility}); hence $\BGA \subsetneq \PS$
(\Cref{cor:strict}), and the generalization is strict.

The conceptual message is that the dual architecture usually attributed to
residuation---dual monoids, De Morgan laws, order reversal---is in fact
generated by three compatibility axioms, \eqref{ax:PJ}, \eqref{ax:TJ1}, and
\eqref{ax:TJ2}, which are jointly much weaker than the residuation law
\eqref{ax:BGA3}.

Several directions remain open. On the axiomatic side, one may ask for the
exact strength of residuation over the PS axioms: which additional axioms,
formulated in the PS signature, characterize the Polar Semirings that do
extend to bounded Girard algebras? On the model-theoretic side, the variety
generated by Polar Semirings invites study of its free objects, its
subdirectly irreducible members, and the decidability of its equational
theory. On the structural side, the non-commutative case, in which
\eqref{ax:T2} is dropped and the polarity may split into left and right
polarities, and the distributive case, in which the induced lattice is
required to be distributive, both appear tractable with the methods of
\Cref{sec:arch}.
\noteAI{If you have partial results on any of these directions (e.g.\ the
characterization of residuable Polar Semirings), consider replacing the
corresponding sentence with a precise conjecture.}
